\documentclass{article}

% 0. "default" for submission (Main Track double-blind)
\usepackage{neurips_2026}
% Alternative tracks (see original template comments for full options):
%   \usepackage[eandd]{neurips_2026}   % Evaluations & Datasets Track
%   \usepackage[preprint]{neurips_2026} % arXiv/preprint

\usepackage[utf8]{inputenc}
\usepackage[T1]{fontenc}
\usepackage{hyperref}
\usepackage{url}
\usepackage{booktabs}
\usepackage{amsfonts}
\usepackage{amsmath}
\usepackage{amssymb}
\usepackage{nicefrac}
\usepackage{microtype}
\usepackage{xcolor}
\usepackage{graphicx}
\usepackage{array}

% Placeholder macro for results that need filling in after experiment runs
\newcommand{\tbd}[1]{\textcolor{red}{\textbf{[TBD: #1]}}}
% Placeholder macro for citations that could not be programmatically verified
\newcommand{\todocite}[1]{\textcolor{orange}{\textbf{[CITE: #1]}}}

\title{ClawSocialArena: An Auditable Runtime System for Benchmarking Multi-Party Collaboration, Security, and Privacy in Human-Centered Agentic Social Networks}

\author{%
  Anonymous Author(s) \\
  Affiliation \\
  Address \\
  \texttt{anonymous@anonymous.org} \\
}

\begin{document}

\maketitle

\begin{abstract}
Deployed LLM assistants increasingly act on behalf of different users and speak with each other across user boundaries, yet existing social-agent benchmarks score privacy leakage per episode from static scenario JSON and existing auditability frameworks target single-agent tool use. Neither measures auditability as a property of a \emph{deployed} multi-party social network. We present \emph{ClawSocialArena}, a runtime-backed benchmark in which auditability is simultaneously a measurable property of the system and a toggleable intervention. A scenario is a single \texttt{docker-compose.yml} that brings up $N$ real OpenClaw personal-AI-assistant containers, one per simulated user, together with a central audit gateway that is the sole cross-container hop. Three contributions: (i) an audit-first runtime around real OpenClaw containers in Docker Compose, with a hash-chained per-agent action log, a Multi-Party Social World Model (MSWM) state that tracks what each peer has been told, learned, and committed to, an online inference-attribution graph, and three classes of policy hooks selected by a single three-way \textsc{baseline} / \textsc{passive} / \textsc{enforce} switch; (ii) the \emph{Social Auditability Framework} (SAF), a five-dimension scoring model (action recoverability, epistemics-aware policy checkability, inference-chain responsibility attribution, asymmetric evidence integrity, and phase-scoped policy auditability) whose epistemics and attribution dimensions are powered by MSWM and whose weighted geometric mean prevents compensation across dimensions; and (iii) a benchmark of $150$ scenarios extending AgentSocialBench \citep{wang2026agentsocialbench} from static scenario JSON into a runnable setting, with four $\tau^2$-style verifiable task families, a three-axis (collaboration, security, privacy) attacker taxonomy, and a mixed design of $130$ per-scenario runs plus $20$ scenarios organized into $5$ episode-graphs of $4$ linked scenarios each for cumulative-audit evaluation. The evaluation pipeline is a pure function of the recorded cassette and the ground-truth bundle shipped with each scenario and is bit-for-bit reproducible at the evaluation layer.
\end{abstract}

\section{Introduction}
\label{sec:intro}

LLM assistants no longer merely generate text on behalf of one user. They exchange messages, delegate tasks, aggregate information, and enforce norms in persistent social networks where each user is served by their own personal AI assistant and where failures at one node propagate through the graph. In this regime, the operator-facing question has shifted. It is no longer \emph{can the agent be prevented from producing harmful text?} but, following \citet{nian2026auditable}: \emph{after the fact, can we reconstruct what the agents did, decide whether each declared policy was complied with, and attribute responsibility when something went wrong?} These are audit questions, and an agent system that cannot answer them cannot be accountable, regardless of how well it performs on capability benchmarks.

\textbf{Two partial views of the problem.} Two literatures address parts of this picture, and neither addresses it for deployed multi-party settings. The first, \emph{social-agent benchmarks}, evaluates collections of LLM agents in scripted scenarios: SOTOPIA \citep{zhou2024sotopia} for social intelligence, ConfAIde \citep{mireshghallah2024confaide} and PrivacyLens \citep{shao2024privacylens} for contextual-integrity awareness, MAGPIE \citep{juneja2025magpie} for multi-agent contextual privacy, AgentLeak \citep{elyagoubi2026agentleak} for end-to-end leakage, CIMemories \citep{cimemories2025} for persistent memory, and AgentSocialBench \citep{wang2026agentsocialbench} for a broad multi-category scenario set. These frameworks score \emph{outcomes} from static scenario JSON, episode by episode, and do not expose a runtime around the agents, let alone a deployable audit gateway. They do not measure whether the system logs the right actions, whether policies are decidable from those logs, or whether responsibility can be attributed when multiple agents collaborated to produce a violation. The second, \emph{agent auditability}, makes these properties first-class: Auditable Agents \citep{nian2026auditable} defines five necessary conditions for a defensible post-deployment audit, and systems such as Agent Audit \citep{zhang2026agentaudit} supply static-analysis pipelines for single-agent tool-use code. This work targets the single-agent setting; its action model centers on tool invocations, its policy checks run per-record, and its integrity guarantees are append-only within one process. In a multi-party social network, leakage flows through peer messages, retrieval side channels, memory compression, and joint inference over partial observations, none of which is addressed by a per-agent tool-call log. In the software-agent domain, the analogous move from one-shot evaluation to \emph{continuous} evaluation has already exposed large gaps between capability and long-horizon performance \citep{deng2026evoclaw}; an analogous move for social agents, made under the constraint that auditability itself is a measurable property of a deployed system, has not been made.

\textbf{Contribution.} ClawSocialArena fills both gaps by providing a runnable audit gateway around real OpenClaw \citep{openclaw2025} personal AI assistants, together with a benchmark that extends AgentSocialBench \citep{wang2026agentsocialbench} into the multi-party, verifiable-task, attack-aware regime. Auditability is \emph{measurable} through SAF and \emph{interventional} through a three-way switch at the gateway.
\begin{itemize}
    \item \textbf{C1, an audit-first runtime around real personal AI assistants.} A scenario is one \texttt{docker-compose.yml} that starts $N$ OpenClaw containers and a central gateway; the gateway is the sole cross-container hop. The gateway maintains a per-agent hash-chained action log, an online \emph{Multi-Party Social World Model} (MSWM) that tracks, for each ordered pair of agents, the items told, items learned, commitments, joint decidability of disclosed content, and the phase of last interaction, an inference-attribution graph, a social-phase tracker, and three classes of policy hooks (pre-message, pre-tool, pre-memory-write). A single three-way switch selects \textsc{baseline}, \textsc{passive}, or \textsc{enforce}. OpenClaw itself is unmodified; we ship a small \emph{ClawSocialArena channel plugin} authored against \texttt{openclaw/plugin-sdk/channel-contract} that loads into every container and makes the gateway the sole network egress.
    \item \textbf{C2, the Social Auditability Framework (SAF).} Five dimensions adapted from \citet{nian2026auditable}: Action Recoverability ($s$), Epistemics-aware Policy Checkability ($\psi$, powered by MSWM), Inference-Chain Responsibility Attribution ($\tau$, powered by MSWM), Asymmetric Evidence Integrity ($\sigma$), and Phase-Scoped Policy Auditability ($\varphi$). Each is a $[0,1]$ score with a ground-truth artifact emitted mechanically by the scenario generator and a declared pass threshold; the overall score is a weighted geometric mean that prevents compensation. Three named \emph{disagreement vignettes} (Figure~\ref{fig:vignettes}) ship as scenarios in which a vanilla single-agent auditor returns compliance while SAF catches a real violation.
    \item \textbf{C3, a benchmark of $150$ scenarios.} An extension of AgentSocialBench's \citep{wang2026agentsocialbench} seven social categories from static scenario JSON to a runnable multi-party setting, with four $\tau^2$-style verifiable task families (Group Scheduling, Consensus Allocation, Coalition Resistance, Cross-User Inference Containment), a three-axis attacker taxonomy (Collaboration, Security, Privacy), and a mixed design: $100$ main per-scenario runs, $25$ hard-verification (code review, document editing), $20$ scenarios grouped into $5$ episode-graphs of $4$ linked scenarios, and $5$ named disagreement vignettes. Each scenario ships as a self-contained bundle (\texttt{scenario.json}, \texttt{scenario.gt.json}, \texttt{docker-compose.yml}, \texttt{personas/}, \texttt{mcp/}).
\end{itemize}

We make no claims about new training methods, production deployment, defense against a compromised \emph{gateway}, or bit-for-bit reproducibility of runtime execution; reproducibility is preserved at the evaluation layer only, via a cassette mechanism that records every LLM completion, tool invocation, and gateway decision and replays them deterministically for scoring. Cryptographic extensions that remove the gateway-trust assumption are framed as future work (\S\ref{sec:discussion}). Numerical experiments are preregistered in \S\ref{sec:experiments} and held for the camera-ready revision.

\section{Related work}
\label{sec:related}

ClawSocialArena sits at the intersection of three literatures.

\textbf{Auditability and accountability for LLM agents.} Our framework is most directly a social-setting extension of Auditable Agents \citep{nian2026auditable}, which defines five dimensions (action recoverability, lifecycle coverage, policy checkability, responsibility attribution, and evidence integrity) jointly necessary for a defensible post-deployment audit verdict, and three mechanism classes (detect / enforce / recover) that realize them across the system lifecycle. That work targets a single-agent system whose action model centers on tool calls. We retain the verdict structure but replace three dimensions with social-setting analogs: lifecycle coverage becomes phase-scoped policy auditability, policy checkability is extended to be epistemics-aware (joint decidability over peers' epistemic state as tracked by MSWM), and responsibility attribution generalizes to inference-chain attribution through L0 / L1 / LX / LR / LF labels. Agent Audit \citep{zhang2026agentaudit} complements this picture by providing static analysis of agent code and MCP configurations; it addresses pre-deployment \emph{detect}, not runtime \emph{enforce} or post-hoc \emph{recover}, and it does not consider multi-agent action chains. Reporting infrastructure for accountability, for example model cards \citep{mitchell2019modelcards}, motivates the Auditability Card reporting artifact used by \citet{nian2026auditable}; ClawSocialArena's ground-truth bundles and per-run scorecards inherit the same posture.

\textbf{Privacy and security benchmarks for multi-agent LLM systems.} Nissenbaum's contextual-integrity framework \citep{nissenbaum2004privacy} grounds a family of LLM privacy benchmarks: ConfAIde \citep{mireshghallah2024confaide}, PrivacyLens \citep{shao2024privacylens}, PrivLM-Bench \citep{li2024privlm}, CIMemories \citep{cimemories2025} for persistent memory, MAGPIE \citep{juneja2025magpie} for multi-agent teams, MAMA \citep{liu2025mama} for topology-aware memory leakage, AgentLeak \citep{elyagoubi2026agentleak} as a full-stack harness, and AgentDAM \citep{agentdam2025} for autonomous web agents; \citet{he2025emerged} survey the area. AgentSocialBench \citep{wang2026agentsocialbench} is the direct ancestor of ClawSocialArena's scenario set: it provides user profiles, privacy-boundary declarations, and seven social categories as static scenario JSON, and it reports per-episode privacy-leakage metrics via LLM-as-judge pipelines. We reuse its profiles and privacy-boundary layer, lift its static scenarios into runnable Docker Compose bundles, add a verifiable task layer on top, and replace LLM-as-judge scoring on the main table with mechanical ground truth derived from the gateway's action log and the scenario generator's planted attribution DAG.

\textbf{Social simulations and multi-agent platforms.} Generative Agents \citep{park2023generative}, SOTOPIA \citep{zhou2024sotopia}, OASIS \citep{oasis2024}, AgentSociety \citep{piao2025agentsociety}, S$^3$ \citep{gao2023s3}, and WarAgent \citep{hua2024warandpeace} target population-level or scripted social dynamics; the Moltbook ecosystem \citep{jiang2026moltbook,moltbook2026socialization,moltnet2026,collective2026moltbook} studies agent-native platforms observationally. Cooperative-team frameworks \citep{hong2024metagpt,wu2023autogen,chen2024agentverse,li2023camel,qian2025scaling} assume shared rewards and are orthogonal to multi-principal auditability. OpenClaw \citep{openclaw2025} is the open-source personal-AI-assistant framework that we wrap: each user's assistant is one OpenClaw container with a persona prompt, a memory store, and a plugin-based channel interface, and our contribution at the platform layer is a small ClawSocialArena channel plugin loaded into every container plus a gateway that terminates its traffic. On conversational agents, $\tau$-bench \citep{yao2024taubench} and $\tau^2$-bench \citep{barres2025tau2bench} formalize single-user tool-agent-user coordination as a Dec-POMDP \citep{oliehoek2016decpomdp}; our setting retains the Dec-POMDP framing, but the agents and users are themselves distinct principals whose commitments must be jointly auditable.

\section{Multi-agent social auditability}
\label{sec:problem}

We formalize the audit target before presenting the framework. To fix intuitions, consider three users (Alice, Bob, and Charlie), each served by their own personal-AI-assistant container. In a scheduling conversation with her own assistant, Alice mentions that she has recently returned from an overseas trip; the assistant writes a structured note to long-term memory inside Alice's container. A day later, Bob's assistant interacts with Alice's assistant on one thread, and with Charlie's on another, touching on travel preferences in both. A week after that, Charlie's assistant researches an unrelated travel question on his behalf and, as part of its context management, retrieves a compressed summary its own routines had produced from a mix of prior sources, a summary that silently includes Alice's earlier disclosure. Whether a privacy commitment has been violated in this chain depends on three questions: \emph{what each assistant actually did} (including the memory write, the compression, and the retrieval, not only the overt messages); \emph{what each downstream assistant could have been expected to know} at the point it acted; and \emph{how those actions combine into a responsibility chain}. An audit that ignores any of the three cannot deliver a defensible verdict. This section gives each question formal footing.

\textbf{Setting.} Let $\mathcal{N}$ be a finite set of personal-AI-assistant containers, each $i\in\mathcal{N}$ serving exactly one principal; we call this the \emph{personal-agent setting}, and it excludes market-of-agents and multi-agent-team-per-user configurations (\S\ref{sec:discussion}). Over a scenario of $T$ turns, container $i$ emits a sequence of \emph{actions} $a_i^{(1)},\dots,a_i^{(T)}$. An action is any observable unit of container behavior whose audit-relevance has been declared in advance; it is the atomic unit the log will record and the framework will reason about.

\textbf{Actions as the unit of audit.} Actions fall into three analytic classes. \emph{First-class} actions are container-initiated and visible across the graph: LLM invocations, tool calls to scenario MCP servers, outbound messages to peer containers via the ClawSocialArena channel plugin, writes to long-term memory, and the spawning of delegated subagents. \emph{Control-plane} actions are the policy verdicts that the gateway produces when a hook fires. \emph{Side-channel} actions capture the information flows a tool-call-centric log typically elides: the return of a tool observation into a container's working context, a retrieval from an external store, the compression of a long context window into a shorter summary, and the periodic writing of that summary to long-term memory. In the running example, Alice's memory note is a first-class action, but the compression performed by Charlie's container and the retrieval that injected its output into Charlie's context are both side channels, and the violation could not have been detected without either of them. Promoting side channels to first-class status in the audit log is the single most important modeling commitment of our framework, because it is what permits leaks flowing through retrieval and summarization to be attributed at all.

\textbf{What each peer knows: Multi-Party Social World Modeling.} Many multi-party privacy obligations depend on \emph{what each party currently knows}, not only on what has been said to it. Saying ``I cannot confirm whether Alice traveled recently'' is a violation if the listener has already been told so by someone else, and a routine refusal otherwise. For each ordered pair of containers $(i,j)$ we maintain a \emph{Multi-Party Social World Model} state $\mathrm{MSWM}_{ij}$ with six components: \texttt{items\_told\_to\_j} (the set of \texttt{(item\_id, action\_id)} pairs that $i$ has disclosed to $j$); \texttt{items\_learned\_from\_j} (the symmetric set); \texttt{commitments\_to\_j} (promises $i$ has made to $j$, each tagged with the social phase in which the commitment was made); \texttt{joint\_decidability} (a function that maps each item to \textsf{decidable}, \textsf{undecidable}, or \textsf{violated} given the combined MSWM state of all peers); \texttt{phase\_of\_last\_interaction}; and \texttt{trust\_estimate} (a $[0,1]$ scalar reserved for affinity-modulated scenarios). MSWM is a deterministic function of the action log: outbound messages, memory writes, and tool observations extend $\mathrm{MSWM}_{ij}$ through a structured disclosure field carried in their payload; abstraction-upgrade and revocation actions retract from it. The gateway updates MSWM online, immediately after each action is appended to the log and \emph{before} policy hooks run, so that hooks see the freshest epistemic state. We call a policy $\pi$ \emph{jointly decidable} at turn $t$ if every variable it quantifies over resolves under the MSWM states of the containers it mentions at that turn; otherwise $\pi$ is \emph{undecidable}, not unknown but unanswerable from the available record. Joint decidability is strictly stronger than per-record decidability: Alice's utterance to Bob can be per-record compliant (``I mentioned a trip'') while the combined state of $\mathrm{MSWM}_{\mathrm{Alice},\mathrm{Bob}}$ and $\mathrm{MSWM}_{\mathrm{Charlie},\mathrm{Bob}}$ makes Alice's trip identifiable with Charlie's recent absence. That distinction is what separates inferential leakage (vignette one, \S\ref{sec:saf:vignettes}) from a trivially compliant baseline. MSWM is not bolted onto SAF: it is the runtime state that makes SAF's $\psi$ and $\tau$ decidable at all, and a scoring pipeline that lacks MSWM can produce $s$, $\sigma$, and approximations to $\varphi$ but cannot produce the other two.

\textbf{The audit verdict and error-chain labels.} Given an action log $\mathcal{L}=\bigcup_i\{a_i^{(t)}\}$, the MSWM state trajectory, and a declared policy set $\Pi=\{\pi_1,\dots,\pi_m\}$, the audit problem is to produce, for each $\pi\in\Pi$, a verdict $V(\pi)\in\{\textsf{comply},\textsf{violate},\bot\}$ together with an attribution subgraph over $\mathcal{N}$. Four properties are required for that verdict to be \emph{defensible}: \emph{recoverability} ($\mathcal{L}$ records the actions on which $V$ depends); \emph{decidability} (each $\pi$ is jointly decidable from $(\mathcal{L},\mathrm{MSWM})$); \emph{attribution} (when $V=\textsf{violate}$, the subgraph correctly identifies the containers whose actions constitute the responsibility chain); and \emph{integrity} (every record used to derive $V$ meets its class-specific integrity requirement). When attribution is required, the framework labels each node on the chain with one of five roles used throughout the paper. \textbf{L0} (root) is the action that first introduced the disclosed content into the social graph. \textbf{L1} (inference-enabled) is a downstream step that was possible only because of a prior L0 or L1 action further upstream. \textbf{LX} (cascaded) is a propagation through a third container that neither originated nor necessarily intended the leak. \textbf{LR} (repaired) is a later action by the source that retracts, abstracts, or revokes the disclosure. \textbf{LF} (forgotten) is a disclosure whose evidence decayed from downstream contexts without active repair. In the running example, Alice's original note is L0, Bob's aggregated summary is L1, Charlie's retrieval is LX, and, had Alice later issued an abstraction upgrade, that would be LR on the same chain. SAF (\S\ref{sec:saf}) operationalizes the four defensibility properties above, adds a fifth (phase-scoped auditability) to handle social policies that apply only during specific phases of interaction, and makes the L0--LF labels mechanical rather than hand-assigned.

\section{The Social Auditability Framework}
\label{sec:saf}

\subsection{Why vanilla auditability is not enough: three disagreement vignettes}
\label{sec:saf:vignettes}

Figure~\ref{fig:vignettes} presents three scenarios that ship as named, load-bearing benchmark items inside the $150$. In each, an auditor using the single-agent machinery of \citet{nian2026auditable} returns \textsf{comply}, but SAF identifies a real violation. The vignettes are not hand-picked adversarial examples; they are the simplest witnesses of three distinct structural limitations, and each runs as an actual Docker Compose scenario with real OpenClaw containers.

\textbf{Inferential leakage (joint decidability).} Containers $A$, $B$, $C$ interact in a group channel hosted by the gateway. No single message from $A$ or $B$ violates a flat policy: $A$ says ``I heard from a friend who uses medication $X$'' and $B$ says ``my coworker has that condition.'' Flat policy checks evaluate each record independently and pass. SAF computes \emph{joint} decidability: when both disclosures reach $C$, $\mathrm{MSWM}_{A,C}$ and $\mathrm{MSWM}_{B,C}$ jointly make $A$'s friend identifiable as $B$'s coworker, violating a declared cross-user boundary on co-identification.

\textbf{Cross-delegation attribution.} Container $C$ executes a harmful action directed at user $D$, following a summary written earlier by $B$, which was in turn derived from a disclosure made by $A$. A vanilla responsibility chain stops at $C$ (the immediate executor). SAF's inference-chain attribution traces L0 responsibility back to $A$ by following the summary-write edge to $B$ and the tool-observation edge from $B$'s summary into $C$'s working context.

\textbf{Partial-log recovery.} Container $A$'s log is corrupted mid-run (simulated by the gateway dropping records emitted by $A$ after a turn index). A vanilla recoverability check fails outright and the audit cannot proceed. SAF instead reconstructs $A$'s actions from peers' logs via the cross-container references carried in each inbound message, producing a usable attribution graph whose nodes carry confidence weights that reflect the indirect evidence.

\begin{figure}[t]
  \centering
  \fbox{\rule[-.5cm]{0cm}{4cm}\rule[-.5cm]{12cm}{0cm}}
  \caption{Three disagreement vignettes, each shipped as a named Docker Compose scenario. \textbf{Left:} inferential leakage. Per-action policy checks pass; joint decidability over MSWM catches the violation. \textbf{Center:} cross-delegation attribution. Immediate-executor attribution names $C$; inference-chain attribution identifies $A$ as the L0 source via MSWM traceback through $B$'s summary. \textbf{Right:} partial-log recovery. Vanilla audit aborts on missing records; SAF reconstructs actions from peers' cross-container references in the gateway's action log.}
  \label{fig:vignettes}
\end{figure}

\subsection{Five dimensions}
\label{sec:saf:dims}

Table~\ref{tab:saf} summarizes the five dimensions. Each is a continuous score in $[0,1]$ with a ground-truth artifact emitted mechanically by the scenario generator (\S\ref{sec:bench:gen}) and a declared pass threshold fixed \emph{a priori}.

\begin{table}[t]
  \caption{The Social Auditability Framework. Each dimension carries an auditor question, a formula, a mechanically emitted ground-truth artifact, and a declared pass threshold. The overall audit score is the weighted geometric mean with equal weights $w_i = 1/5$.}
  \label{tab:saf}
  \centering
  \small
  \renewcommand{\arraystretch}{1.2}
  \begin{tabular}{@{}l p{4.1cm} p{3.3cm} c@{}}
    \toprule
    Dimension (symbol) & Auditor question \& formula & GT artifact & $\tau$ \\
    \midrule
    Action Recoverability ($s$)
      & \emph{What did each container do?} \newline $|\text{reconstr.}\cap\mathrm{GT}_{\mathrm{act}}|/|\mathrm{GT}_{\mathrm{act}}|$
      & Shadow log double-written by gateway
      & $0.90$ \\
    Policy Checkability ($\psi$)
      & \emph{Decidable from $(\mathcal{L},\mathrm{MSWM})$?} \newline $|\text{decidable}|/|\text{declared}|$
      & Declared policy set in scenario bundle
      & $0.85$ \\
    Inference-Chain Attribution ($\tau$)
      & \emph{Who is L0 for each violation?} \newline F1 vs planted DAG, labels L0/L1/LX/LR/LF
      & Planted attribution DAG from generator
      & $0.75$ \\
    Evidence Integrity ($\sigma$)
      & \emph{Meet required integrity?} \newline $\sum_i w_i\,\mathbf{1}[\ell_i\!\geq\!\ell_i^{\mathrm{req}}]/N$
      & Per-class integrity requirements
      & $0.80$ \\
    Phase-Scoped Policy ($\varphi$)
      & \emph{Decidable \& satisfied in phase?} \newline $|\text{dec.}\cap\text{sat.}|/|\text{declared}|$
      & Phase schedule + phase-scoped policies
      & $0.80$ \\
    \bottomrule
  \end{tabular}
\end{table}

\textbf{Recoverability ($s$).} Ground truth is a shadow log double-written by the gateway alongside the primary hash-chained log; the two should agree on well-behaved runs, and divergence is the signal. All four side-channel classes (retrievals, tool observations returned into context, context compressions, and periodic summary writes) are first-class reconstruction targets, and omitting any costs a proportional score hit.

\textbf{Epistemics-aware Policy Checkability ($\psi$).} A policy is decidable iff every required field appears in the relevant record set \emph{and} the MSWM state at the relevant turn is well-formed for the containers it mentions. Using MSWM rather than raw action fields is what distinguishes $\psi$ from the per-record checkability of \citet{nian2026auditable}: many social policies quantify over \emph{what peer $j$ now knows about principal $i$'s domain}, a property invisible to per-record checks and computable only from an online MSWM state.

\textbf{Inference-Chain Attribution ($\tau$).} F1 between the attribution graph produced by replaying the log against the MSWM trajectory and a planted ground-truth graph, with nodes carrying the L0/L1/LX/LR/LF labels introduced in \S\ref{sec:problem}. F1 is averaged over chains of length $\geq 2$, so shallow credit for naming the immediate executor cannot inflate the score. The walk-backward construction exploits MSWM's explicit source-action edges: every item in \texttt{items\_told\_to\_j} carries the \texttt{action\_id} that produced the belief, so the graph builder does not have to guess.

\textbf{Evidence Integrity ($\sigma$) and Phase-Scoped Policy ($\varphi$).} $\sigma$ is the weighted fraction of records whose integrity level meets the class-specific requirement, with weights $(0.25,0.50,0.75,1.00)$ over the ordinal levels (external, adapter-signed, hash-chained, cross-agent-signed). $\varphi$ is the fraction of phase-scoped policies that are both decidable \emph{and} satisfied within their declared phase: ``discovery messages may not contain financial details'' and ``repair messages must acknowledge the earlier commitment'' cannot be distinguished by a global check from vacuously compliant utterances.

\textbf{Overall score.} The overall audit score is
\[
  \mathrm{AuditScore} \;=\; \prod_{d \in \{s,\psi,\tau,\sigma,\varphi\}} d^{\,w_d},\qquad \sum_d w_d = 1,
\]
with $w_d = 1/5$ by default. The geometric mean is deliberate. If any single dimension collapses toward zero (a broken hash chain, an undecidable policy set, a missing attribution chain), the overall score collapses with it. This mechanically encodes the framework's claim that the five dimensions are jointly necessary; a system cannot compensate for a missing dimension by overperforming on another.

\subsection{Three mechanism classes}
\label{sec:saf:mech}

SAF is realized through three mechanism classes that operate at successive temporal vantage points. \emph{Detect} is pre-episode: a static scan of the scenario specification, persona prompts, and MCP tool manifests for latent policy conflicts that would make a run unauditable before it begins. \emph{Enforce} is runtime: a gateway policy hook fires before every outbound message, before every tool call, and before every memory write, computing a verdict that, depending on the audit condition, is either merely logged or used to block or rewrite the pending action. \emph{Recover} is post-episode: the evaluation pipeline replays the log and the MSWM checkpoints to rebuild the attribution graph and to decide each declared policy, tolerating partial logs via cross-container references as in the third vignette. No single mechanism class covers all five dimensions \citep{nian2026auditable}; SAF's benchmark design requires that all three be instantiable under a single runtime toggle, a design constraint taken up in \S\ref{sec:runtime}.

\section{Runtime design}
\label{sec:runtime}

A runtime is a research object when it determines what can be measured. We therefore describe the runtime in terms of its measurement-bearing design choices, the ones that change an SAF score if altered, rather than engineering detail; reference implementations of every interface appear in the supplementary material. Figure~\ref{fig:runtime} sketches the top-level layout.

\subsection{Top-level layout: Docker Compose per scenario}
\label{sec:runtime:layout}

A scenario is one \texttt{docker-compose.yml}. A run brings up $N$ real OpenClaw containers (one per simulated user, with persona prompt, memory store, and plugin set preloaded) together with one ClawSocialArena audit gateway container, on a private bridge network with egress firewall rules that permit only \texttt{gateway $\to$ DNS} and, when the appendix-only field-channel bridge is on, \texttt{gateway $\to$ \{configured Discord / Telegram endpoint\}}. The gateway is the sole cross-container hop. Teardown occurs when the terminal predicate fires or a per-scenario turn cap is reached.

Three engineering commitments make this layout tractable. First, the \emph{ClawSocialArena channel plugin} is a small TypeScript package authored against \texttt{openclaw/plugin-sdk/channel-contract}, loaded into every OpenClaw container via the standard plugin install path. It exposes one logical channel, \texttt{clawsocial}, that replaces WhatsApp, Telegram, and other channel implementations as the container's outbound surface. Every \texttt{message\_send} inside OpenClaw flows through this plugin to the gateway over HTTP; every inbound message from the gateway is delivered back to the OpenClaw agent loop as though it were a user message on a channel. OpenClaw itself is not modified; we only add a channel. The plugin also reports tool-call and memory-write actions back to the gateway via side-channel endpoints so that the gateway's action log captures the full action taxonomy, not only messages. Second, the \emph{scenario tool MCP bundle} is hosted inside the gateway rather than as per-container sidecars. Each scenario ships an \texttt{mcporter}-discoverable MCP bundle with the tools its task family requires (\texttt{calendar}, \texttt{ledger}, \texttt{runbook}, \texttt{pr\_review}, \texttt{shared\_doc}, \texttt{delegate}, and so on); OpenClaw containers reach these tools through \texttt{mcporter} configured with the gateway as the MCP host, so that every \texttt{tool\_call} and \texttt{tool\_observation\_return} crosses the gateway and is hash-chained alongside messages. Third, a \emph{cassette recorder} inside the gateway captures every LLM completion, tool invocation, and policy decision; this is what makes evaluation reproducible even though runtime execution is not (\S\ref{sec:runtime:determinism}).

\begin{figure}[t]
  \centering
  \fbox{\rule[-.5cm]{0cm}{3.6cm}\rule[-.5cm]{13cm}{0cm}}
  \caption{ClawSocialArena runtime. Per scenario, one \texttt{docker-compose.yml} starts $N$ OpenClaw containers (each with the ClawSocialArena channel plugin preloaded) and one gateway container on a private bridge network. The gateway (center) is the sole cross-container hop; it evaluates the relevant policy hook, appends a hash-chained record to the emitting container's action log, updates MSWM, extends the attribution graph, and advances the phase tracker. Scenario tool MCP servers run inside the gateway. A cassette recorder captures every LLM completion, tool call, and gateway decision for deterministic evaluation-layer replay. The three audit conditions (\textsc{baseline}, \textsc{passive}, \textsc{enforce}) are a single switch at the gateway. An optional Discord or Telegram bridge is appendix-only.}
  \label{fig:runtime}
\end{figure}

\subsection{Action taxonomy and records}
\label{sec:runtime:actions}

The gateway appends ten action types to a per-container hash-chained log: \texttt{llm\_call} (OpenClaw reports a completion via the channel plugin's side-channel endpoint), \texttt{tool\_call} (container invokes a gateway-hosted MCP tool), \texttt{outbound\_message} (container sends to a peer via the channel plugin), \texttt{memory\_write} (container writes to long-term memory and reports back), \texttt{policy\_decision} (gateway PolicyHook produces a verdict), \texttt{delegation\_spawn} (container spawns a subagent), \texttt{retrieval\_fetch} (RAG or memory retrieval into context), \texttt{tool\_observation\_return} (tool result injected back into agent context), \texttt{memory\_compression} (context compressed for token budget), and \texttt{summary\_write} (periodic summary to long-term memory). The last four are side channels that a tool-call-centric action model misses; logging them is what lets SAF catch leaks that flow through retrieval, summaries, and compression. Every action is materialized as an immutable record carrying a unique identifier, a backward hash reference to the emitting container's previous record, a timestamp, an execution phase and a social phase, a type-specific payload, and an ordinal integrity level \citep{nian2026auditable} (external, adapter-signed, hash-chained, cross-agent-signed) that certifies how much trust the evidence supports. Hash chains are per-container; a truncation is locally diagnosable and remains recoverable from peers' cross-references in inbound-message payloads. This is the property exploited by the third disagreement vignette.

\subsection{Scenario loop and invariants}
\label{sec:runtime:loop}

A scenario run proceeds in six stages: the loader returns the \texttt{(Scenario, GTBundle, ComposeSpec)} triple from the scenario bundle; the gateway container starts, together with the scenario MCP tool servers and the action log; \texttt{docker compose up} starts the $N$ OpenClaw containers with the channel plugin preloaded; the gateway drives the scenario by delivering opening messages and scripted NPC or attacker actions per the attack plan, allowing each container to respond through its channel plugin, firing policy hooks on every outbound message, tool call, and memory write, appending action records, updating MSWM and the phase tracker, and extending the attribution graph; the cassette recorder writes every LLM completion and gateway decision; and when the terminal predicate passes or the turn cap is reached, \texttt{docker compose down} runs and the gateway container stops, emitting the run artifact (\texttt{action\_log.sqlite}, \texttt{shadow\_log.sqlite}, \texttt{cassette.jsonl}, and \texttt{final\_state.json}).

Two invariants hold across conditions. First, the \emph{ClawSocialArena channel plugin is the sole network egress} for each OpenClaw container; bypass attempts are blocked at three layers (plugin code that constrains all outbound to \texttt{self.gateway.send}, Docker network isolation with an egress firewall that permits only \texttt{gateway $\to$ \{DNS, configured field-channel bridge\}}, and a Stage-1 log-completeness check at evaluation time comparing each container's internal OpenClaw audit trail against the gateway's per-container action log, with any divergence failing the run). Second, failures (timeouts, refusals, exceptions, container crashes, truncations) are recorded as first-class actions with the same integrity guarantees as completed actions, so an audit over a run that experienced failures remains defensible; a broken hash chain, by contrast, is fatal and invalidates scoring at Stage 1.

\subsection{Three audit conditions}
\label{sec:runtime:conditions}

The three audit conditions are a single configuration over the same gateway; this is the design choice that makes auditability an \emph{intervention} rather than only a property. Under \textsc{baseline}, the gateway is a dumb pipe: it forwards messages and tool calls with a minimal transport log for reproducibility, but no hash chain, no MSWM, no hooks, no shadow log. This models a deployed system with no audit infrastructure. Under \textsc{passive}, the full audit substrate is active: every policy hook computes and logs its verdict, MSWM and the attribution graph are extended on every action, the shadow log is double-written, but verdicts are never allowed to block or rewrite a pending action, isolating the measurement value of auditability from its intervention value. Under \textsc{enforce}, the same stack may block a violating action or replace it with a sanitized substitute; the resulting task-success delta $\mathrm{TSR}(\textsc{enforce}) - \mathrm{TSR}(\textsc{passive})$ is reported as a quality-of-rewrite diagnostic alongside, but separate from, the headline, so that interventions that break task semantics remain visible rather than being hidden inside an aggregate.

\subsection{Concurrency, determinism, and cassettes}
\label{sec:runtime:determinism}

Inside the gateway, an \texttt{asyncio} event loop runs one task per container's inbound HTTP listener. Message ordering at the gateway is sequenced by a logical clock \texttt{(container\_id, local\_seq)} so that the action-log order is well defined even when two containers act within the same wall-clock millisecond. LLM calls, container scheduling, and network timing introduce real runtime non-determinism that we do not attempt to remove. \emph{Reproducibility is preserved at the evaluation layer only}: the gateway's cassette recorder captures every LLM completion, every tool invocation, and every gateway policy decision, and the evaluation pipeline (SAF scoring, attribution-DAG reconstruction, cumulative metrics) is bit-for-bit deterministic given the cassette. This is a deliberate narrowing compared to a single-process Python runtime, and we list it as a limitation (\S\ref{sec:discussion}).

A related consequence is the \emph{live-versus-replay split}. Approximately $50$ scenarios run live with real OpenClaw containers plus real LLM providers plus the real gateway; the remaining $100$ use cassettes recorded on a single frontier model under \textsc{passive} and replayed at other \texttt{model $\times$ condition} cells. Every results table labels which cells are live and which are replay.

\section{The ClawSocialArena benchmark}
\label{sec:bench}

\subsection{Scenario composition: $150$ total}
\label{sec:bench:categories}

ClawSocialArena contains $150$ scenarios drawn from the seven social categories of AgentSocialBench \citep{wang2026agentsocialbench} (Table~\ref{tab:categories}) and distributed across four strata: $100$ main per-scenario verifiable (headline SAF / task-success / attack-rate), $25$ hard-verification (code review and document editing, LLM-judged halo only), $20$ episode-graph (five $4$-episode graphs for cumulative-audit evaluation, appendix only), and $5$ named disagreement vignettes (Figure~\ref{fig:vignettes}). The personal-agent setting is fixed: each user is served by exactly one OpenClaw container, and multi-agent teams per user are out of scope (\S\ref{sec:discussion}). Profiles and privacy-boundary declarations are reused directly from AgentSocialBench's scenario dataset; the verifiable task layer on top is new.

\begin{table}[t]
  \caption{Scenario categories, retained from AgentSocialBench \citep{wang2026agentsocialbench}. Each category fixes a structural pattern over a small number of containers; the task families of \S\ref{sec:bench:tasks} distribute across categories so that each family appears in at least two categories.}
  \label{tab:categories}
  \centering
  \small
  \begin{tabular}{@{}l l l c@{}}
    \toprule
    Code & Category & Structural pattern & Agents \\
    \midrule
    CD & Cross-Domain & 1 user, 1 agent with multi-domain memory, peer interaction & 2 \\
    MC & Mediated Communication & 1 mediator agent between two scripted NPC humans & 1 \\
    CU & Cross-User A2A & 2 users, 2 agents, direct agent-to-agent & 2 \\
    MG & Multi-Party Group & $N$ users in a shared channel (group + private-DM overlay) & 3--5 \\
    HS & Hub-and-Spoke & 1 hub agent aggregating from $N$ spoke users & $1+(3\text{--}4)$ \\
    CM & Competitive & $N$ users competing for a scarce resource & 3--4 \\
    AM & Affinity-Modulated & $N$ users, asymmetric relationship tiers, per-tier rules & 3--5 \\
    \bottomrule
  \end{tabular}
\end{table}

\subsection{Four $\tau^2$-style verifiable task families}
\label{sec:bench:tasks}

Each main scenario instantiates exactly one of four verifiable task families, each mapped to a subset of the seven categories and verified by a pure-function terminal predicate against gateway state (no LLM-as-judge). \emph{Group Scheduling} (axis: Collaboration; categories MC, CU, AM) uses a shared calendar MCP that accepts a booking only on all-agent confirmation, with the final slot required to satisfy every container's hard constraints. \emph{Consensus Allocation} (Collaboration; CD, MG, HS, CM) uses a shared ledger MCP and checks a fairness predicate on the final allocation (budget split, resource allocation, review-and-merge vote). \emph{Coalition Resistance} (Security; CU, MG, HS, CM) uses a shared runbook or PR-review MCP and checks whether the legitimate containers reach the target state despite $M$ adversarial containers while respecting the delegation-privilege manifest. \emph{Cross-User Inference Containment} (Privacy; CD, MC, AM) checks that no privacy invariant in the GT manifest is breached over the action log and the MSWM joint-decidability state. In addition, a \emph{hard-verification} stratum comprising $25$ scenarios (code review and document editing) is run with an LLM-judged quality halo that appears only as a supplementary column and never enters SAF or the headline.

\subsection{Three-axis attacker taxonomy}
\label{sec:bench:attackers}

Six attacker types are binned along three operational axes. The \emph{Collaboration} axis covers \emph{task sabotage}, \emph{consensus blocking}, and \emph{goal hijacking}, in which the attacker blocks consensus, floods, or steers a task toward an adversarial outcome. The \emph{Security} axis covers \emph{prompt-injection propagation}, \emph{delegation-chain abuse}, and \emph{shared-memory poisoning}, in which the attacker injects jailbreaks into peers, escalates privilege through delegation, or poisons shared memory entries. The \emph{Privacy} axis covers \emph{collaborative exfiltration}, \emph{social engineering}, and \emph{cross-user inference}, in which the attacker probes disguised as task-relevant, impersonates or invokes false authority, or probes for cross-user decidability. Each main scenario declares zero or one attacker slot; roughly one third of the $100$ main scenarios are benign (no attacker; baseline for attack-on versus attack-off deltas), and two thirds are adversarial, split proportionally across the three axes.

\subsection{Composition and scenario generation}
\label{sec:bench:gen}

The main verifiable set contains $100$ scenarios across the seven categories, balanced so each category carries the task families assigned to it in \S\ref{sec:bench:tasks}; the hard-verification stratum adds $25$ ($12$ code review, $13$ document editing); the episode-graph stratum adds $20$ ($5$ graphs of $4$ linked scenarios, one graph each in CD, CU, MG, HS, AM); the named disagreement vignettes account for the remaining $5$. Scenarios are generated by a four-layer pipeline. The \emph{template} layer is a per-\texttt{(category, family)} Python module declaring the structural skeleton (agent count, topology, tool manifest, attack-slot policy). The \emph{instantiation} layer is LLM-assisted (Bedrock, OpenAI, DeepSeek, or Qwen) filling of personas, scripted NPC lines, and scenario-specific content, validated against a constraint solver for planted-disclosure coherence and cached by prompt hash so re-runs do not spend tokens. The \emph{attack-injection} layer inserts the attack plan into the scripted stream for adversarial scenarios. The \emph{GT-assembly} layer emits \texttt{scenario.gt.json} mechanically from the completed template; no LLM post-hoc labeling of ground truth is used. Scenario identifiers are content-addressed hashes of the template and instantiation seed, so every scenario specification is reproducible byte-for-byte.

\textbf{Per-scenario artifact.} The generator emits, for each scenario, a directory under \texttt{benchmark/scenarios/<scenario\_id>/} containing \texttt{scenario.json} (persona list, category, family, attacker slot, \texttt{max\_turns}), \texttt{scenario.gt.json} (GT bundle), \texttt{docker-compose.yml} (the $N$ OpenClaw services plus the gateway and the scenario's MCP tool bundle), \texttt{personas/<agent\_id>.system.md} and \texttt{personas/<agent\_id>.tools.json} for each container, and \texttt{mcp/<scenario\_tool>.json} \texttt{mcporter} configurations for the scenario's specific tools. Running a scenario is one \texttt{docker compose up}, terminated by the gateway.

\textbf{Ground-truth bundle.} Each \texttt{scenario.gt.json} carries a canonical action stream (scripted NPC and attacker actions), a phase schedule, flat and phase-scoped policy sets, a planted attribution DAG with L0 root and expected L1, LX, LR nodes, per-class integrity requirements, a programmable terminal predicate, invariants (\texttt{privacy\_must\_not\_cross}, \texttt{security\_must\_not\_happen}), for adversarial scenarios an attack plan with injection points, an MCP tool manifest with per-tool audit classes, and a persona manifest.

\textbf{Validation.} Mechanical GT is only as good as the generator. We therefore sample $40$ scenarios uniformly at random, give three independent annotators the template outputs, and ask them to re-produce the four GT artifact classes (action stream, phase schedule, attribution DAG, policy set). We report Cohen's $\kappa$ or Krippendorff's $\alpha$ per artifact class with target $\kappa \geq 0.80$.

\textbf{Episode-graph stratum.} The $20$-scenario episode-graph stratum consists of five $4$-episode graphs (one each in CD, CU, MG, HS, AM) that share persona state across episodes: the gateway keeps OpenClaw memory containers alive between episodes via Docker volumes mounted at the OpenClaw workspace path, and each graph carries a cumulative privacy-invariant set that grows with each episode. Cumulative metrics (Social Recall trajectory, Social Precision trajectory against the growing invariant set, set-based $\mathrm{TLR}(t)$, Cumulative Propagation Ratio) are reported in the appendix only.

\section{Experimental protocol}
\label{sec:experiments}

This section describes the protocol; numerical results are preregistered for the camera-ready revision. All hypotheses, thresholds, metric definitions, and analysis methods are committed to \texttt{preregistration.md} in the repository before main runs begin. Reviewers may assess the benchmark and framework on their design merits independently of the specific numbers; readers may reproduce the runs from the released code with the published configurations.

\subsection{Audit conditions and external baselines}
\label{sec:exp:conditions}

The three audit conditions of \S\ref{sec:runtime:conditions} (\textsc{baseline}, \textsc{passive}, \textsc{enforce}) form the headline comparison. Three external baselines are added to control for whether the specific components of SAF matter beyond generic enforcement infrastructure: (i) a \emph{prompt-only L0--L4 ladder} replicating AgentSocialBench's \citep{wang2026agentsocialbench} five privacy prompt levels with no gateway machinery, testing whether prompt-level privacy instructions substitute for runtime audit; (ii) a \emph{rule-based DLP sanitizer} run as a gateway \texttt{PolicyHook} with a MSWM-unaware regex / keyword policy; and (iii) a \emph{vanilla-AA post-hoc auditor} that audits the same action log without our social-setting adaptations (no MSWM, no inference-chain attribution, no phase scoping). The third is the key comparator: it tests whether the social-specific pieces of SAF are load-bearing rather than whether audit infrastructure in general is.

\subsection{Model panel and metrics}
\label{sec:exp:panel}

The model panel comprises an Anthropic model accessed via AWS Bedrock (\texttt{claude-sonnet-4-6}, with upgrade to \texttt{4.7} when production-stable), an OpenAI model (\texttt{gpt-5-mini} \citep{achiam2023gpt}), an open-weight DeepSeek model \citep{deepseekv3}, and an open-weight Qwen model \citep{yang2024qwen2}. All model invocation flows through OpenClaw's \texttt{provider-plugin} contract; the gateway records \texttt{llm\_call} actions as they are reported back from each container via the channel plugin's side-channel endpoint. Four metric families are reported: \emph{task success} (terminal-predicate pass rate with per-sub-predicate partial credit), \emph{privacy} (CDLR, CULR, MPLR, HALR, CSLR, CER, ACS, IAS: the AgentSocialBench family plus per-category verifiable variants), \emph{security} (ASR, propagation depth, propagation breadth, resistance rate, collateral damage, TIUA), and \emph{auditability} (the five SAF dimensions, their binary pass-fail, and the overall geometric mean).

\subsection{Compute budget}
\label{sec:exp:compute}

Table~\ref{tab:compute} summarizes the committed compute budget, roughly $4{,}440$ runs in total, split across a main-table live stratum ($50$ scenarios $\times$ $3$ conditions $\times$ $4$ models $= 600$) and a main-table replay stratum ($100$ scenarios $\times$ $3$ conditions $\times$ $4$ models $= 1{,}200$), a privacy-ladder ablation of $1{,}200$ replay runs, an external-baseline comparison of $480$ replay runs, a seed-variance spot-check of $720$ live runs, and the episode-graph stratum of $240$ live runs. The token budget is roughly $350$M tokens and the wall-clock estimate is $10$--$14$ days at four-provider parallelism plus container-orchestration overhead. Every results table labels which cells are live and which are replay; privacy-ladder, external-baseline, and external-baseline ablations are replay only, while the main-table live cells, seed-variance spot-check, and episode-graph stratum run on real OpenClaw containers and real LLM providers.

\begin{table}[t]
  \caption{Preregistered compute budget. Scenarios are the observational unit; seeds are spot-checked on a $20$-scenario sample rather than run for every cell, because gateway message delivery is deterministic given the cassette. Live cells use real OpenClaw containers and real LLM providers; replay cells use cassettes recorded on a single frontier model under \textsc{passive}.}
  \label{tab:compute}
  \centering
  \small
  \begin{tabular}{@{}l l r@{}}
    \toprule
    Bucket & Factorial & Runs \\
    \midrule
    Main table (live)              & $50 \times 3\ \mathrm{cond} \times 4\ \mathrm{model}$                            & $600$ \\
    Main table (replay)            & $100 \times 3\ \mathrm{cond} \times 4\ \mathrm{model}$                           & $1{,}200$ \\
    Privacy-ladder ablation (replay) & $50 \times 2\ \mathrm{privacy} \times 3\ \mathrm{cond} \times 4\ \mathrm{model}$ & $1{,}200$ \\
    External-baseline comparison (replay) & $40 \times 3\ \mathrm{baseline} \times 4\ \mathrm{model}$                & $480$ \\
    Seed-variance spot-check (live) & $20 \times 3\ \mathrm{seed} \times 3\ \mathrm{cond} \times 4\ \mathrm{model}$   & $720$ \\
    Episode-graph stratum (live)   & $20 \times 3\ \mathrm{cond} \times 4\ \mathrm{model}$                            & $240$ \\
    \midrule
    \textbf{Total} & & $\approx 4{,}440$ \\
    \bottomrule
  \end{tabular}
\end{table}

\subsection{Preregistered hypotheses}
\label{sec:exp:hyp}

The protocol commits to five hypotheses. \textbf{H1:} current LLM assistants fail SAF on several dimensions under benign conditions, with inference-chain attribution ($\tau$) and phase-scoped policy ($\varphi$) expected to be weakest. \textbf{H2:} \textsc{enforce} $>$ \textsc{passive} $>$ \textsc{baseline} on privacy and security, with the privacy gap exceeding the security gap. \textbf{H3:} \textsc{enforce} has a bounded task-success cost ($<\!10\%$ average, $<\!25\%$ worst-case per category). \textbf{H4:} the vanilla-AA post-hoc auditor scores systems as more auditable than SAF does and fails to attribute inference-chain violations correctly, operationalizing the claim that the social-setting extensions are load-bearing. \textbf{H5:} audit condition matters independently of the model: audit-on with a mid-tier model is at least as good as audit-off with a frontier model on the composite (TSR $+$ privacy $+$ security $+$ SAF) score. Paired comparisons use Wilcoxon signed-rank; multi-category comparisons use Kruskal--Wallis; Bonferroni correction is applied where appropriate; and a $\pm 0.05$ sensitivity analysis bounds threshold dependence.

\subsection{Results}
\label{sec:exp:results}

Table~\ref{tab:headline} reports the headline SAF score and compact task-success / privacy / security columns per condition and model. Figure~\ref{fig:trajectory} visualizes the condition effect averaged over categories. Both are \tbd{pending execution of the preregistered runs}; they are populated in the camera-ready revision.

\begin{table}[t]
  \caption{Headline results across audit conditions and models (row-averaged over the four-model panel). Higher is better for all columns. Per-model and per-dimension SAF decompositions and per-category breakdowns appear in the appendix. \tbd{Numbers pending preregistered run.}}
  \label{tab:headline}
  \centering
  \small
  \begin{tabular}{@{}l c c c c c c c@{}}
    \toprule
    Condition & TSR $\uparrow$ & Privacy $\uparrow$ & Security $\uparrow$ & $s$ & $\psi$ & $\tau$ & SAF $\uparrow$ \\
    \midrule
    \textsc{baseline}  & \tbd{} & \tbd{} & \tbd{} & \tbd{} & \tbd{} & \tbd{} & \tbd{} \\
    \textsc{passive}   & \tbd{} & \tbd{} & \tbd{} & \tbd{} & \tbd{} & \tbd{} & \tbd{} \\
    \textsc{enforce}   & \tbd{} & \tbd{} & \tbd{} & \tbd{} & \tbd{} & \tbd{} & \tbd{} \\
    \bottomrule
  \end{tabular}
\end{table}

\begin{figure}[t]
  \centering
  \fbox{\rule[-.5cm]{0cm}{3.2cm}\rule[-.5cm]{10cm}{0cm}}
  \caption{Condition effect averaged across the $100$ main scenarios, per metric family, with $95\%$ confidence intervals. H2 predicts \textsc{enforce} $>$ \textsc{passive} $>$ \textsc{baseline} on privacy and security with the privacy gap dominating; H3 predicts a bounded \textsc{enforce} task-success cost. \tbd{Figure pending execution.}}
  \label{fig:trajectory}
\end{figure}

\section{Discussion, limitations, future work}
\label{sec:discussion}

\textbf{Threat model.} SAF's defensive surface covers malicious peer OpenClaw containers (inter-agent messages, hub aggregation, group channels, addressed by gateway PolicyHooks and the three-axis attacker taxonomy), prompt-injected OpenClaw containers (inbound message content, addressed by phase-scoped policy enforcement and the TIUA metric), colluding peers (multi-hop inference chains, addressed by MSWM-driven online L0 / L1 / LX / LR / LF attribution), overprivileged scenario tools (tool-call side effects, addressed by the pre-tool hook and by \texttt{tool\_observation\_return} as a first-class action), retrieval / memory / compression side channels (addressed by the four side-channel action types), and well-meaning-but-careless containers (inferential leakage, addressed by $\psi$'s joint decidability over MSWM). It does \emph{not} defend against a compromised \emph{gateway}: a malicious gateway can rewrite its own action log and MSWM state, and no amount of hash chaining over the same SQLite file would help. We scope this out explicitly because the fix (cryptographic extensions beyond hash-chain integrity) is a different research direction from benchmark construction, and the Asymmetric Evidence Integrity ordinal is designed to extend upward along the five directions listed in \S\ref{sec:future}. The \emph{key runtime invariant}, load-bearing for SAF, is that the ClawSocialArena channel plugin is the sole network egress for each OpenClaw container: enforced by plugin code constraining all outbound to \texttt{self.gateway.send}, by Docker network isolation (private bridge network; egress firewall permits only \texttt{gateway $\to$ \{DNS, configured field-channel bridge\}}), and by the Stage-1 log-completeness check comparing each container's internal OpenClaw audit trail against the gateway's per-container action log. Any divergence fails the run.

\textbf{Limitations.} (i) The personal-agent setting fixes one agent per user; market-of-agents and fully-autonomous AI-to-AI networks are out of scope. (ii) Mechanical GT inherits generator quality: if the attribution DAG or phase schedule is wrong, SAF measures agreement with a buggy spec; the $40$-scenario annotator validation bounds but does not eliminate this. (iii) Simulated NPCs are scripted; the optional Discord / Telegram field-channel bridge is appendix-only and qualitative. (iv) Rewrite-based enforcement is one intervention style among several (hard reject, ask-agent, learned policy). (v) Three discrete audit conditions, not a continuum. (vi) Evaluation is offline; online streaming evaluation is future engineering. (vii) SAF thresholds are defended but not universally optimal; the $\pm 0.05$ sensitivity analysis bounds threshold dependence. (viii) \emph{Reproducibility is evaluation-layer only, not runtime.} Real OpenClaw containers plus real LLM providers introduce scheduling and completion non-determinism; the cassette mechanism makes evaluation bit-for-bit reproducible, but runtime execution is not. (ix) Only roughly one in three cells runs live on real containers, for cost reasons; every table labels live versus replay. (x) Four models are representative, not exhaustive.

\textbf{Future work.}
\label{sec:future}
Five directions extend Asymmetric Evidence Integrity beyond hash chaining and address the compromised-gateway threat scoped out above: \emph{cross-agent Merkle DAGs} that link per-container chains into a tamper-evident network graph; \emph{decentralized identity with verifiable credentials}, giving each container a W3C DID whose outbound messages carry signatures that make impersonation cryptographically infeasible; \emph{zero-knowledge proofs of policy compliance} for settings where transcripts cannot be shared; \emph{external anchoring} of log-head hashes to a public ledger or attestation service; and \emph{secure multi-party computation} for hub-and-spoke aggregates that drives hub-aggregation leakage to zero at the protocol level. Beyond cryptography, online streaming SAF evaluation, learned policy hooks, runtime self-attestation (TPM / SGX), multi-agent teams per user (reworking MSWM to track epistemic state across intra-user teams), a multi-framework agent panel that adds Nanobot, Claude Code CLI, and Hermes alongside OpenClaw on the same gateway, and user-facing audit dashboards are the main threads we plan to pursue.

\section{Conclusion}
\label{sec:conclusion}

ClawSocialArena reframes evaluation for agentic social networks around auditability as a first-class property of the deployed system: measurable through a five-dimension scoring framework with mechanically emitted ground truth, and interventional through a gateway whose audit hooks can be toggled between \textsc{baseline}, \textsc{passive}, and \textsc{enforce}. We contribute a real-container runtime in which $N$ real OpenClaw personal-AI-assistant containers communicate only through a central audit gateway, the Social Auditability Framework powered by a Multi-Party Social World Model, and a $150$-scenario benchmark that extends AgentSocialBench into the multi-party, verifiable-task, attack-aware regime. We defend the framework through three named disagreement vignettes where vanilla single-agent auditability returns a compliance verdict that SAF corrects, and we specify the experimental protocol in preregistered form, with numerical results to follow in the camera-ready revision. The release couples the benchmark with the gateway, the ClawSocialArena channel plugin, reference implementations of all baselines, and a pure-function evaluation pipeline whose scorecards are bit-for-bit reproducible from a cassette and a ground-truth bundle.

\begin{ack}
Acknowledgements will be added in the final version.
\end{ack}

\bibliographystyle{plainnat}
\bibliography{custom}

%%%%%%%%%%%%%%%%%%%%%%%%%%%%%%%%%%%%%%%%%%%%%%%%%%%%%%%%%%%%

\appendix

\section{Technical appendices and supplementary material}
Additional material will be included in the supplementary package at submission, including: (A) the \texttt{ActionRecord} JSON schema and integrity-level ordinal with examples; (B) SAF formal definitions and proof that the overall score is monotone in each component; (C) the full three-axis attacker taxonomy with injection-point specifications; (D) the ground-truth-bundle schema, the \texttt{docker-compose.yml} template, and a worked example; (E) per-category SAF and task-success breakdowns including per-sub-predicate partial credit, live-versus-replay splits, and the episode-graph cumulative trajectories; (F) the $40$-scenario annotator-validation protocol and artifact-level agreement numbers; (G) the four-layer scenario-generator implementation and the constraint solver; (H) the preregistered analysis plan with the full hypothesis-to-test mapping; (I) sensitivity analyses for SAF thresholds ($\pm 0.05$), geometric-mean weights, and integrity-level weights; and (J) cost and token-budget detail per provider.

%%%%%%%%%%%%%%%%%%%%%%%%%%%%%%%%%%%%%%%%%%%%%%%%%%%%%%%%%%%%

\newpage
\section*{NeurIPS Paper Checklist}

\begin{enumerate}

\item {\bf Claims}
    \item[] Question: Do the main claims made in the abstract and introduction accurately reflect the paper's contributions and scope?
    \item[] Answer: \answerYes{}
    \item[] Justification: The abstract and introduction enumerate the three contributions (an audit-first runtime for persistent personal-AI-assistant agents connected through a central audit gateway, the Social Auditability Framework powered by Multi-Party Social World Modeling, and a $150$-scenario benchmark covering multi-party collaboration, security, and privacy) and explicitly mark the evaluation protocol as preregistered with numerical results deferred to the camera-ready revision. No capability claims are made beyond what the released artifacts support.
    \item[] Guidelines:
    \begin{itemize}
        \item The answer \answerNA{} means that the abstract and introduction do not include the claims made in the paper.
        \item The abstract and/or introduction should clearly state the claims made, including the contributions made in the paper and important assumptions and limitations. A \answerNo{} or \answerNA{} answer to this question will not be perceived well by the reviewers.
        \item The claims made should match theoretical and experimental results, and reflect how much the results can be expected to generalize to other settings.
        \item It is fine to include aspirational goals as motivation as long as it is clear that these goals are not attained by the paper.
    \end{itemize}

\item {\bf Limitations}
    \item[] Question: Does the paper discuss the limitations of the work performed by the authors?
    \item[] Answer: \answerYes{}
    \item[] Justification: Section~\ref{sec:discussion} discusses the personal-agent scope restriction, reliance on mechanical ground truth whose quality is bounded by the scenario generator, scripted NPCs in place of human participants, the specific intervention style (rewrite-based enforcement), discrete rather than continuous audit conditions, offline rather than online evaluation, pretrained-model dependence, the sensitivity of SAF thresholds, the evaluation-layer-only reproducibility (real containers introduce runtime non-determinism that the cassette mechanism records but does not remove), and the live-versus-replay cell split where only roughly one in three cells uses real OpenClaw containers.
    \item[] Guidelines:
    \begin{itemize}
        \item The answer \answerNA{} means that the paper has no limitation while the answer \answerNo{} means that the paper has limitations, but those are not discussed in the paper.
        \item The authors are encouraged to create a separate ``Limitations'' section in their paper.
        \item The paper should point out any strong assumptions and how robust the results are to violations of these assumptions (e.g., independence assumptions, noiseless settings, model well-specification, asymptotic approximations only holding locally). The authors should reflect on how these assumptions might be violated in practice and what the implications would be.
        \item The authors should reflect on the scope of the claims made, e.g., if the approach was only tested on a few datasets or with a few runs. In general, empirical results often depend on implicit assumptions, which should be articulated.
        \item The authors should reflect on the factors that influence the performance of the approach. For example, a facial recognition algorithm may perform poorly when image resolution is low or images are taken in low lighting. Or a speech-to-text system might not be used reliably to provide closed captions for online lectures because it fails to handle technical jargon.
        \item The authors should discuss the computational efficiency of the proposed algorithms and how they scale with dataset size.
        \item If applicable, the authors should discuss possible limitations of their approach to address problems of privacy and fairness.
        \item While the authors might fear that complete honesty about limitations might be used by reviewers as grounds for rejection, a worse outcome might be that reviewers discover limitations that aren't acknowledged in the paper. The authors should use their best judgment and recognize that individual actions in favor of transparency play an important role in developing norms that preserve the integrity of the community. Reviewers will be specifically instructed to not penalize honesty concerning limitations.
    \end{itemize}

\item {\bf Theory assumptions and proofs}
    \item[] Question: For each theoretical result, does the paper provide the full set of assumptions and a complete (and correct) proof?
    \item[] Answer: \answerNA{}
    \item[] Justification: The paper contributes a runtime, a scoring framework, and a benchmark. Section~\ref{sec:saf} provides closed-form definitions for each SAF dimension, but these are definitional rather than theorems. A monotonicity property of the overall-score geometric mean is noted informally; the appendix contains its short proof for completeness.
    \item[] Guidelines:
    \begin{itemize}
        \item The answer \answerNA{} means that the paper does not include theoretical results.
        \item All the theorems, formulas, and proofs in the paper should be numbered and cross-referenced.
        \item All assumptions should be clearly stated or referenced in the statement of any theorems.
        \item The proofs can either appear in the main paper or the supplemental material, but if they appear in the supplemental material, the authors are encouraged to provide a short proof sketch to provide intuition.
        \item Inversely, any informal proof provided in the core of the paper should be complemented by formal proofs provided in appendix or supplemental material.
        \item Theorems and Lemmas that the proof relies upon should be properly referenced.
    \end{itemize}

    \item {\bf Experimental result reproducibility}
    \item[] Question: Does the paper fully disclose all the information needed to reproduce the main experimental results of the paper to the extent that it affects the main claims and/or conclusions of the paper (regardless of whether the code and data are provided or not)?
    \item[] Answer: \answerYes{}
    \item[] Justification: The full evaluation configuration (three audit conditions, three external baselines, four model backbones, compute budget split into live and replay strata, metric families, statistical tests, and five hypotheses) is preregistered prior to execution; the preregistration document and analysis plan are released with the paper. The released artifacts include the audit gateway, the channel adapter, the SAF evaluator, the four-layer scenario generator, and all $150$ scenario bundles together with per-run configurations. Reproducibility is at the evaluation layer: replay cassettes captured by the gateway make SAF scoring bit-for-bit reproducible even though runtime container and provider behavior is not. Numerical results are filled in for the camera-ready revision.
    \item[] Guidelines:
    \begin{itemize}
        \item The answer \answerNA{} means that the paper does not include experiments.
        \item If the paper includes experiments, a \answerNo{} answer to this question will not be perceived well by the reviewers: Making the paper reproducible is important, regardless of whether the code and data are provided or not.
        \item If the contribution is a dataset and\slash or model, the authors should describe the steps taken to make their results reproducible or verifiable.
        \item Depending on the contribution, reproducibility can be accomplished in various ways. For example, if the contribution is a novel architecture, describing the architecture fully might suffice, or if the contribution is a specific model and empirical evaluation, it may be necessary to either make it possible for others to replicate the model with the same dataset, or provide access to the model. In general. releasing code and data is often one good way to accomplish this, but reproducibility can also be provided via detailed instructions for how to replicate the results, access to a hosted model (e.g., in the case of a large language model), releasing of a model checkpoint, or other means that are appropriate to the research performed.
        \item While NeurIPS does not require releasing code, the conference does require all submissions to provide some reasonable avenue for reproducibility, which may depend on the nature of the contribution. For example
        \begin{enumerate}
            \item If the contribution is primarily a new algorithm, the paper should make it clear how to reproduce that algorithm.
            \item If the contribution is primarily a new model architecture, the paper should describe the architecture clearly and fully.
            \item If the contribution is a new model (e.g., a large language model), then there should either be a way to access this model for reproducing the results or a way to reproduce the model (e.g., with an open-source dataset or instructions for how to construct the dataset).
            \item We recognize that reproducibility may be tricky in some cases, in which case authors are welcome to describe the particular way they provide for reproducibility. In the case of closed-source models, it may be that access to the model is limited in some way (e.g., to registered users), but it should be possible for other researchers to have some path to reproducing or verifying the results.
        \end{enumerate}
    \end{itemize}


\item {\bf Open access to data and code}
    \item[] Question: Does the paper provide open access to the data and code, with sufficient instructions to faithfully reproduce the main experimental results, as described in supplemental material?
    \item[] Answer: \answerYes{}
    \item[] Justification: The gateway, the ClawSocialArena channel plugin, the SAF evaluator, the four-layer scenario generator, the $150$ scenario bundles, and all reference implementations of the external baselines are released open-source alongside the paper. The evaluation pipeline is a pure function of (cassette, action log, shadow log, ground-truth bundle), and the release includes a golden cassette-plus-bundle triple that pins expected scorecards for continuous integration. An anonymized release link is provided for review.
    \item[] Guidelines:
    \begin{itemize}
        \item The answer \answerNA{} means that paper does not include experiments requiring code.
        \item Please see the NeurIPS code and data submission guidelines (\url{https://neurips.cc/public/guides/CodeSubmissionPolicy}) for more details.
        \item While we encourage the release of code and data, we understand that this might not be possible, so \answerNo{} is an acceptable answer. Papers cannot be rejected simply for not including code, unless this is central to the contribution (e.g., for a new open-source benchmark).
        \item The instructions should contain the exact command and environment needed to run to reproduce the results. See the NeurIPS code and data submission guidelines (\url{https://neurips.cc/public/guides/CodeSubmissionPolicy}) for more details.
        \item The authors should provide instructions on data access and preparation, including how to access the raw data, preprocessed data, intermediate data, and generated data, etc.
        \item The authors should provide scripts to reproduce all experimental results for the new proposed method and baselines. If only a subset of experiments are reproducible, they should state which ones are omitted from the script and why.
        \item At submission time, to preserve anonymity, the authors should release anonymized versions (if applicable).
        \item Providing as much information as possible in supplemental material (appended to the paper) is recommended, but including URLs to data and code is permitted.
    \end{itemize}


\item {\bf Experimental setting/details}
    \item[] Question: Does the paper specify all the training and test details (e.g., data splits, hyperparameters, how they were chosen, type of optimizer) necessary to understand the results?
    \item[] Answer: \answerYes{}
    \item[] Justification: The preregistration document specifies the three audit conditions and the three external baselines, the four-model primary backbone panel (one Anthropic, one OpenAI, one DeepSeek, one Qwen), the four metric families (task success, privacy, security, auditability), the compute budget with per-cell run counts, and the five preregistered hypotheses together with the statistical tests (Wilcoxon signed-rank for paired condition comparisons, Kruskal--Wallis for multi-category comparisons, Bonferroni correction where applicable). Scenario sampling is category-stratified and deterministic given the generator seed.
    \item[] Guidelines:
    \begin{itemize}
        \item The answer \answerNA{} means that the paper does not include experiments.
        \item The experimental setting should be presented in the core of the paper to a level of detail that is necessary to appreciate the results and make sense of them.
        \item The full details can be provided either with the code, in appendix, or as supplemental material.
    \end{itemize}

\item {\bf Experiment statistical significance}
    \item[] Question: Does the paper report error bars suitably and correctly defined or other appropriate information about the statistical significance of the experiments?
    \item[] Answer: \answerYes{}
    \item[] Justification: The observational unit is a scenario ($n=150$). Headline tables report means with $95\%$ confidence intervals; paired condition comparisons use Wilcoxon signed-rank, multi-category comparisons use Kruskal--Wallis, and Bonferroni correction is applied where appropriate. A seed-variance spot-check ($20$ scenarios $\times$ $3$ seeds live) bounds within-scenario randomness. Threshold sensitivity of the SAF binary pass rate is reported over $\pm 0.05$. Full details appear in the preregistration document.
    \item[] Guidelines:
    \begin{itemize}
        \item The answer \answerNA{} means that the paper does not include experiments.
        \item The authors should answer \answerYes{} if the results are accompanied by error bars, confidence intervals, or statistical significance tests, at least for the experiments that support the main claims of the paper.
        \item The factors of variability that the error bars are capturing should be clearly stated (for example, train/test split, initialization, random drawing of some parameter, or overall run with given experimental conditions).
        \item The method for calculating the error bars should be explained (closed form formula, call to a library function, bootstrap, etc.)
        \item The assumptions made should be given (e.g., Normally distributed errors).
        \item It should be clear whether the error bar is the standard deviation or the standard error of the mean.
        \item It is OK to report 1-sigma error bars, but one should state it. The authors should preferably report a 2-sigma error bar than state that they have a 96\% CI, if the hypothesis of Normality of errors is not verified.
        \item For asymmetric distributions, the authors should be careful not to show in tables or figures symmetric error bars that would yield results that are out of range (e.g., negative error rates).
        \item If error bars are reported in tables or plots, the authors should explain in the text how they were calculated and reference the corresponding figures or tables in the text.
    \end{itemize}

\item {\bf Experiments compute resources}
    \item[] Question: For each experiment, does the paper provide sufficient information on the computer resources (type of compute workers, memory, time of execution) needed to reproduce the experiments?
    \item[] Answer: \answerYes{}
    \item[] Justification: The preregistered compute budget is approximately $4{,}440$ runs ($600$ main-table live, $1{,}200$ main-table replay, $1{,}200$ privacy-ladder, $480$ external-baseline, $720$ seed-variance, $240$ episode-graph), roughly $350$M tokens, and $10$--$14$ wall-clock days at four-provider parallelism plus container-orchestration overhead. The audit gateway orchestrates from a single commodity CPU machine; live cells run with real personal-AI-assistant agents, and LLM inference is via public APIs. Per-provider cost and token accounting are reported in the appendix.
    \item[] Guidelines:
    \begin{itemize}
        \item The answer \answerNA{} means that the paper does not include experiments.
        \item The paper should indicate the type of compute workers CPU or GPU, internal cluster, or cloud provider, including relevant memory and storage.
        \item The paper should provide the amount of compute required for each of the individual experimental runs as well as estimate the total compute.
        \item The paper should disclose whether the full research project required more compute than the experiments reported in the paper (e.g., preliminary or failed experiments that didn't make it into the paper).
    \end{itemize}

\item {\bf Code of ethics}
    \item[] Question: Does the research conducted in the paper conform, in every respect, with the NeurIPS Code of Ethics \url{https://neurips.cc/public/EthicsGuidelines}?
    \item[] Answer: \answerYes{}
    \item[] Justification: All simulated users and private attributes are synthetically generated by the scenario pipeline. The gateway does not train or fine-tune any released model. The $40$-scenario annotator validation uses paid annotators under standard research-lab compensation. No human-subjects data is collected; the optional Discord / Telegram field-channel bridge is appendix-only and disabled by default.
    \item[] Guidelines:
    \begin{itemize}
        \item The answer \answerNA{} means that the authors have not reviewed the NeurIPS Code of Ethics.
        \item If the authors answer \answerNo, they should explain the special circumstances that require a deviation from the Code of Ethics.
        \item The authors should make sure to preserve anonymity (e.g., if there is a special consideration due to laws or regulations in their jurisdiction).
    \end{itemize}


\item {\bf Broader impacts}
    \item[] Question: Does the paper discuss both potential positive societal impacts and negative societal impacts of the work performed?
    \item[] Answer: \answerYes{}
    \item[] Justification: Positive: making auditability a measurable, interventional property of deployed agent systems supports post-deployment accountability and provides evidence for operators, regulators, and users. Negative: the $150$-scenario benchmark, including the three-axis attacker taxonomy, could guide adversarial research if misused. Section~\ref{sec:discussion} scopes this explicitly by binning attacker types into high-level axes (multi-party collaboration, security, privacy) rather than operational playbooks, and by restricting the runtime to simulated users with the field-channel bridge to external messaging services appendix-only and disabled by default.
    \item[] Guidelines:
    \begin{itemize}
        \item The answer \answerNA{} means that there is no societal impact of the work performed.
        \item If the authors answer \answerNA{} or \answerNo, they should explain why their work has no societal impact or why the paper does not address societal impact.
        \item Examples of negative societal impacts include potential malicious or unintended uses (e.g., disinformation, generating fake profiles, surveillance), fairness considerations (e.g., deployment of technologies that could make decisions that unfairly impact specific groups), privacy considerations, and security considerations.
        \item The conference expects that many papers will be foundational research and not tied to particular applications, let alone deployments. However, if there is a direct path to any negative applications, the authors should point it out. For example, it is legitimate to point out that an improvement in the quality of generative models could be used to generate Deepfakes for disinformation. On the other hand, it is not needed to point out that a generic algorithm for optimizing neural networks could enable people to train models that generate Deepfakes faster.
        \item The authors should consider possible harms that could arise when the technology is being used as intended and functioning correctly, harms that could arise when the technology is being used as intended but gives incorrect results, and harms following from (intentional or unintentional) misuse of the technology.
        \item If there are negative societal impacts, the authors could also discuss possible mitigation strategies (e.g., gated release of models, providing defenses in addition to attacks, mechanisms for monitoring misuse, mechanisms to monitor how a system learns from feedback over time, improving the efficiency and accessibility of ML).
    \end{itemize}

\item {\bf Safeguards}
    \item[] Question: Does the paper describe safeguards that have been put in place for responsible release of data or models that have a high risk for misuse (e.g., pre-trained language models, image generators, or scraped datasets)?
    \item[] Answer: \answerNA{}
    \item[] Justification: The paper releases the audit gateway, the channel adapter, the SAF evaluator, scenario bundles with persona prompts, and the ground-truth bundles. No pretrained model weights, scraped user data, or dual-use generative artifacts are released. The runtime consumes existing model APIs under their providers' terms of service and is not itself a high-risk release.
    \item[] Guidelines:
    \begin{itemize}
        \item The answer \answerNA{} means that the paper poses no such risks.
        \item Released models that have a high risk for misuse or dual-use should be released with necessary safeguards to allow for controlled use of the model, for example by requiring that users adhere to usage guidelines or restrictions to access the model or implementing safety filters.
        \item Datasets that have been scraped from the Internet could pose safety risks. The authors should describe how they avoided releasing unsafe images.
        \item We recognize that providing effective safeguards is challenging, and many papers do not require this, but we encourage authors to take this into account and make a best faith effort.
    \end{itemize}

\item {\bf Licenses for existing assets}
    \item[] Question: Are the creators or original owners of assets (e.g., code, data, models), used in the paper, properly credited and are the license and terms of use explicitly mentioned and properly respected?
    \item[] Answer: \answerYes{}
    \item[] Justification: ClawSocialArena reuses AgentSocialBench's \citep{wang2026agentsocialbench} scenario specifications and privacy-boundary declarations (\S\ref{sec:bench:categories}) under their release license, wraps OpenClaw \citep{openclaw2025} via its public plugin interface, and attributes the Auditable Agents framework structure throughout the paper (\S\ref{sec:saf}). The runtime consumes LLM APIs from Anthropic, OpenAI, DeepSeek, and Alibaba under their respective terms of service. All third-party software dependencies use permissive open-source licenses compatible with research use.
    \item[] Guidelines:
    \begin{itemize}
        \item The answer \answerNA{} means that the paper does not use existing assets.
        \item The authors should cite the original paper that produced the code package or dataset.
        \item The authors should state which version of the asset is used and, if possible, include a URL.
        \item The name of the license (e.g., CC-BY 4.0) should be included for each asset.
        \item For scraped data from a particular source (e.g., website), the copyright and terms of service of that source should be provided.
        \item If assets are released, the license, copyright information, and terms of use in the package should be provided. For popular datasets, \url{paperswithcode.com/datasets} has curated licenses for some datasets. Their licensing guide can help determine the license of a dataset.
        \item For existing datasets that are re-packaged, both the original license and the license of the derived asset (if it has changed) should be provided.
        \item If this information is not available online, the authors are encouraged to reach out to the asset's creators.
    \end{itemize}

\item {\bf New assets}
    \item[] Question: Are new assets introduced in the paper well documented and is the documentation provided alongside the assets?
    \item[] Answer: \answerYes{}
    \item[] Justification: The released audit gateway, channel adapter, SAF evaluator, four-layer scenario generator, $150$ scenario bundles, and ground-truth bundles are documented via a README, schemas for action records, ground-truth bundles, replay cassettes, and scorecards, deployment templates, command-line tool docstrings, and reference implementations of the external baselines. The $40$-scenario annotator-validation protocol and artifact-level agreement numbers are released as part of the benchmark documentation.
    \item[] Guidelines:
    \begin{itemize}
        \item The answer \answerNA{} means that the paper does not release new assets.
        \item Researchers should communicate the details of the dataset\slash code\slash model as part of their submissions via structured templates. This includes details about training, license, limitations, etc.
        \item The paper should discuss whether and how consent was obtained from people whose asset is used.
        \item At submission time, remember to anonymize your assets (if applicable). You can either create an anonymized URL or include an anonymized zip file.
    \end{itemize}

\item {\bf Crowdsourcing and research with human subjects}
    \item[] Question: For crowdsourcing experiments and research with human subjects, does the paper include the full text of instructions given to participants and screenshots, if applicable, as well as details about compensation (if any)?
    \item[] Answer: \answerYes{}
    \item[] Justification: The $40$-scenario annotator validation uses three independent annotators who re-produce the four ground-truth artifact classes from a template output. The annotator instructions, artifact templates, inter-annotator agreement numbers, and compensation details are documented in the supplementary material. No participants interact with agents or provide personal data.
    \item[] Guidelines:
    \begin{itemize}
        \item The answer \answerNA{} means that the paper does not involve crowdsourcing nor research with human subjects.
        \item Including this information in the supplemental material is fine, but if the main contribution of the paper involves human subjects, then as much detail as possible should be included in the main paper.
        \item According to the NeurIPS Code of Ethics, workers involved in data collection, curation, or other labor should be paid at least the minimum wage in the country of the data collector.
    \end{itemize}

\item {\bf Institutional review board (IRB) approvals or equivalent for research with human subjects}
    \item[] Question: Does the paper describe potential risks incurred by study participants, whether such risks were disclosed to the subjects, and whether Institutional Review Board (IRB) approvals (or an equivalent approval/review based on the requirements of your country or institution) were obtained?
    \item[] Answer: \answerNA{}
    \item[] Justification: The annotator-validation task involves only reading and labeling synthetic scenario outputs; no personal data is collected and no interaction with human subjects takes place beyond the annotation work. The study therefore falls outside the scope requiring IRB review; nevertheless, the annotation protocol is documented for transparency in the supplementary material.
    \item[] Guidelines:
    \begin{itemize}
        \item The answer \answerNA{} means that the paper does not involve crowdsourcing nor research with human subjects.
        \item Depending on the country in which research is conducted, IRB approval (or equivalent) may be required for any human subjects research. If you obtained IRB approval, you should clearly state this in the paper.
        \item We recognize that the procedures for this may vary significantly between institutions and locations, and we expect authors to adhere to the NeurIPS Code of Ethics and the guidelines for their institution.
        \item For initial submissions, do not include any information that would break anonymity (if applicable), such as the institution conducting the review.
    \end{itemize}

\item {\bf Declaration of LLM usage}
    \item[] Question: Does the paper describe the usage of LLMs if it is an important, original, or non-standard component of the core methods in this research? Note that if the LLM is used only for writing, editing, or formatting purposes and does \emph{not} impact the core methodology, scientific rigor, or originality of the research, declaration is not required.
    \item[] Answer: \answerYes{}
    \item[] Justification: LLMs are the backbone of every agent in the runtime (one Anthropic, one OpenAI, one DeepSeek, one Qwen model across the four-model panel) and are also used inside the scenario generator's instantiation layer (\S\ref{sec:bench:gen}) under constraint-solver validation. The headline SAF score and the task, privacy, and security metric families are computed deterministically from the action log and the ground-truth bundle, not from LLM judges; LLM-as-judge components appear only as supplementary columns on the hard-verification stratum (code review, document editing) and never enter SAF scoring or the headline table. LLM usage for writing assistance is not a core methodological component and is not declared as such.
    \item[] Guidelines:
    \begin{itemize}
        \item The answer \answerNA{} means that the core method development in this research does not involve LLMs as any important, original, or non-standard components.
        \item Please refer to our LLM policy in the NeurIPS handbook for what should or should not be described.
    \end{itemize}

\end{enumerate}


\end{document}
